%! TeX program = lualatex
% One page: the map of mindmap-limits.tex page 1, turned a quarter turn and
% set one size larger.  It lives in its own file because it is a print
% variant of that map rather than another page of it -- keeping it here means
% adding a page there never renumbers what the notes ask for.
%
% `multi' is what makes the mmpage a page of its own -- the [tikz] option
% would do that for a bare tikzpicture, but it leaves the picture ungrouped
% and \resizebox below then fails with a missing \endgroup.
\documentclass[multi=mmpage]{standalone}
\input{../typography.tex.preamble}

% DO NOT use the tikz preamble: its `basic curve' styles pull in pgfplots, and
% enabling `trig format=rad' there makes this picture die with a dimension too
% large error.
% \input{../tikz.tex.preamble}
\usepackage{graphicx}
\usepackage{tikz}

\usetikzlibrary{mindmap, backgrounds, calc}

% The "substitution just works" case.  WesternSpring on its own is a neon
% lime; muting it through grey keeps it in the palette without shouting.
\colorlet{substcase}{WesternSpring!55!gray!45}
% Same hue at full strength, for the arrow joining two substcase nodes.
\colorlet{substline}{WesternSpring!55!gray}
% Both one-sided limits exist.  WesternSky neat is a neon cyan, so mute it
% through grey the same way.
\colorlet{bothexist}{WesternSky!60!gray!45}
% The named theorems the map tells you to reach for -- the Intermediate
% Value Theorem and Squeeze.  They sit in different branches, so a colour of
% their own is what marks them as the same kind of thing: a theorem you
% invoke, rather than a case in somebody's decision tree.
\colorlet{namedthm}{blue!20}

% Connector labels.  Neutral and italic rather than coloured: the node fills
% already carry the colour coding, and `attn' gold went muddy beside them.
\colorlet{decision}{black!60}

\newenvironment{mmpage}{}{}

% Second line of a concept label.  Formulas stay in \displaystyle, as
% \everymath sets course-wide, so \lim subscripts sit underneath: the circles
% below are sized to allow for that.  One size larger than the same pair in
% mindmap-limits.tex: this is the less wide of the two layouts, so it is
% scaled down less and can afford the extra room.
\newcommand{\subB}[1]{\\[0.3ex]{\small\(#1\)}}
\newcommand{\noteB}[1]{\\[0.2ex]{\footnotesize #1}}

% The mindmap library gives a level-2 concept 1.5cm of text width and a
% level-3 one just 1cm, which hyphenates single words like remov-able and
% Inter-mediate.  Widen them, and keep the manual line breaks short enough to
% fit what is set here.
\tikzset{
  mindmap base/.style={
    mindmap,
    grow cyclic,
    concept color=gray!20,
    every child node/.style={concept, minimum size=1.4cm, text width=1.7cm,
                             align=center, font=\scriptsize},
  },
  % Connector labels.  `mindmap' hands every node a large minimum size; that is
  % harmless when a label is centred on a point, but an edge-anchored one ends
  % up centimetres off its bar unless the box is shrunk back to its text.
  barlabel/.style={font=\footnotesize\itshape, text=decision, align=center,
                   inner sep=2pt, minimum size=0pt,
                   minimum width=0pt, minimum height=0pt},
  barlabelB/.style={barlabel, font=\normalsize\itshape},
}

% Two things about the angles here.
%
% (1) A `clockwise from' given to the root's children leaks down to every
%     deeper level, so without a correction every node's children fan in the
%     same absolute direction instead of pointing outward.  Each node
%     therefore restates its own fan, in the path right after it -- in a
%     child's option list the key would instead re-aim that child itself.
%     The value is (parent angle) - (children - 1) * (sibling angle) / 2.
%
% (2) Level distances and sibling angles differ per subtree, so they are set
%     in each child's option list, where they propagate down that subtree.

\begin{document}

% ===================================================================
%  The same map as mindmap-limits.tex page 1, turned a quarter turn.  The
%  long axis of the layout (continuity on one side, the techniques on the
%  other) then runs vertically, so the drawing is taller than wide and
%  fills the page.
% ===================================================================
\begin{mmpage}
% The guided notes want a 7in-wide graphic (mindmap-integration.pdf is 7.3in).
% The picture is drawn wider than that and scaled down to fit, so keep labels
% short: at this size every extra line of text costs legibility.
\resizebox{7in}{!}{%
\begin{tikzpicture}[
  mindmap base,
  level 1/.append style={level distance=5.77cm, sibling angle=120},
  level 2/.append style={level distance=4.20cm, sibling angle=60},
  level 3/.append style={level distance=3.77cm, sibling angle=44},
  ]
  % Branches at 135, 45, -45, -135.  Continuity sits next to direct
  % substitution so the cross-connection between them stays short.
  % The number line is a nested tikzpicture inside the root's label: no tick
  % labels, one dot at a, and the axis tip named x.
  \node [concept, minimum size=4.51cm, text width=4.12cm, font=\large] (root)
    {limit at a constant \(a\)\\[0.8ex]
     % The nested picture inherits the mindmap's node sizing, which flings
     % these two labels centimetres away; reset it here.
     \begin{tikzpicture}[baseline,
         every node/.style={inner sep=1.5pt, minimum size=0pt, minimum width=0pt,
                            minimum height=0pt, text width=, align=center}]
       \draw [->, thin] (0,0) -- (2.5,0) node [right] {\scriptsize\(x\)};
       \fill (1.25,0) circle (1.4pt);
       \node [below] at (1.25,0) {\scriptsize\(a\)};
     \end{tikzpicture}}
    [clockwise from=240]
    child [concept color=gray!20,
           every node/.append style={text width=2.71cm},
           level 2/.append style={level distance=6.10cm, sibling angle=50},
           level 3/.append style={level distance=3.77cm, sibling angle=52}] { 
      % Wide enough to keep the test on one line -- a node option beats the
      % 2.1cm this subtree sets for everything below it.
      node [minimum size=3.80cm, text width=4.38cm, font=\normalsize] (conttest)
        {continuity at \(a\)\subB{\lim_{x \to a} f(x) \overset{?}{=} f(a)}}
      [counterclockwise from=165]
      child [concept color=substcase] { 
        node (yes) {continuous}
        [counterclockwise from=155]
        child [counterclockwise from=120] { node (families) [text width=3.35cm] {can use direct\\substitution\noteB{to evaluate limits of\\polynomials, rational,\\trigs, \(e^{x}\), \(\ln(x)\)\\in their domain}} }
        child [concept color=namedthm, counterclockwise from=158, level distance=2.70cm] { node (ivt) {Intermediate\\Value Theorem} }
      }
      child { node (leftcont) [minimum size=2.61cm, text width=3.09cm]
        {left-continuous\noteB{\(\lim_{x \to a^{-}} f(x) \overset{?}{=} f(a)\)}} }
      child { node (rightcont) [minimum size=2.61cm, text width=3.09cm]
        {right-continuous\noteB{\(\lim_{x \to a^{+}} f(x) \overset{?}{=} f(a)\)}} }
      child [concept color=gray!20,
             level 3/.append style={level distance=5.45cm, sibling angle=52}] { 
        node (no) [text width=3.20cm] {discontinuous:\noteB{calculate both\\one-sided limits}}
        [counterclockwise from=270]
        % Cyan for the two cases where both one-sided limits exist -- not gold,
        % which is the connector-label colour.  The infinite case keeps the
        % pink it shares with nonzero/0.
        child [concept color=bothexist] { node (removable) {removable} }
        child [concept color=bothexist] { node (jump) {jump} }
        child [concept color=pink!50] { node (infinite) [text width=3.35cm]
          {infinite and \(f\) has\\a vertical\\asymptote at \(a\)} }
      }
    }
    % Pushed further out than its siblings: the bar into it carries a
    % two-line label, and the techniques hanging off it need the room.
    child [concept color=gray!20, level distance=6.95cm,
           level 3/.append style={level distance=4.40cm, sibling angle=40}] { 
      node [minimum size=2.85cm, text width=3.50cm, font=\normalsize] (dirsub)
        {try direct\\substitution\noteB{is the denominator \(0\)?}}
      % Squeeze hangs off `try direct substitution' itself rather than off
      % 0/0.  It goes first, so the fan starts one sibling angle earlier (0
      % rather than 60) and the three cases keep their angles; its own level
      % distance puts it within 0.3cm of where it sat as a grandchild, and
      % `missing' holds its old place there.
      [counterclockwise from=0]
      child [concept color=namedthm, level distance=5cm] { node (squeeze) {Squeeze\\Theorem} }
      child [concept color=supp!20] { 
        node (zerozero) {\(\frac{0}{0}\)}
        % Four children now, and the fan is swung well below this node's own
        % direction so the last of them clears the nonzero/0 subtree above.
        [counterclockwise from=-49]
        child [missing]
        child { node (cov) {change of\\variable} }
        child { node (rationalize) {rationalize} }
        child { node (factor) {factor} }
        % Its own colour, reserved for differentiation elsewhere in the notes.
        child [concept color=main!25] { node (deriv) [text width=2.83cm]
          {recognize as\\a derivative} }
      }
      % Pink to pair it with the infinite discontinuity the arrow points at.
      child [concept color=pink!50,
             level 3/.append style={level distance=5.45cm}] { 
        node (constzero) {\(\frac{\text{nonzero}}{0}\)}
        [counterclockwise from=100]
        child { node (differ) {limit does\\not exist} }
        child { node (agree) {\(+\infty\) or \(-\infty\)} }
      }
      child [concept color=substcase] { node (algno) [text width=2.71cm]
          {\(\frac{\text{any number}}{\text{nonzero}}\)} }
    }
    child [concept color=gray!20,
           level 2/.append style={level distance=5.24cm}] { 
      node (twosided) [minimum size=2.61cm, text width=2.96cm]
        {two-sided limit\subB{\lim_{x \to a} f(x)}}
      [counterclockwise from=-60]
      child { node (leftlim) [minimum size=2.61cm, text width=2.96cm]
        {left limit at \(a\)\subB{\lim_{x \to a^{-}} f(x)}} }
      child { node (rightlim) [minimum size=2.61cm, text width=2.96cm]
        {right limit at \(a\)\subB{\lim_{x \to a^{+}} f(x)}} }
      child [concept color=supp!20] { node (known) [text width=2.58cm]
        {well-known\\limits} }
    }
    ;

  \node [concept color=white, annotation, text width=, above,
         font=\small, align=center] at (known.north)
    % `aligned' rather than `align*': a display environment inside a TikZ node
    % dies with a missing \endgroup.  \everymath is \displaystyle, so this
    % sets exactly as align* would.
    {\(\begin{aligned}
       \lim_{x \to 0} \frac{\sin(x)}{x} &= 1\\
       \lim_{x \to 0} \frac{\cos(x) - 1}{x} &= 0\\
       \lim_{x \to 0} \frac{e^{x} - 1}{x} &= 1\\
       \lim_{x \to \infty} \left(1 + \frac{1}{x}\right)^{x} &= e
     \end{aligned}\)};

  % Below the node, its left edge flush with the node's.  A circle's own
  % .south west sits at 225 degrees, inset from the left edge, so take the
  % x of .west and the y of .south instead.
  \node [concept color=white, annotation, anchor=north west, text width=5.41cm,
         font=\small, align=left] at ([yshift=-2pt]ivt.south -| ivt.west)
    {IVT applies to continuous function over \([a,b]\) and \(f(a) \ne f(b)\)};

  % Labels riding the bars into a branch, rather than sitting in it as one
  % more node.  All of them are decisions or instructions, so they share the
  % `attn' colour and read as a layer distinct from the concept names.
  % In the foreground, not on the arrow's own path: the arrow is drawn on the
  % background layer, so a label riding it disappears under the nodes it passes.
  % Sloped along its own arrow.  An invisible foreground path between the same
  % two nodes carries the label, since the arrow itself is on the background
  % layer and anything riding it would be hidden by the nodes it passes.
  \path (cov) -- (known)
    node [barlabelB, sloped, below=2pt, midway] {in some cases,\\we use \dots};
  \node [barlabelB] at ($(no)!0.5!(removable)$)
    {both exist and agree,\\but \(\ne f(a)\)};
  \node [barlabelB] at ($(no)!0.5!(jump)$) {both exist\\but disagree};
  \node [barlabelB] at ($(no)!0.5!(infinite)$)
    {one of them\\is \(\pm\infty\)};
  \node [barlabelB] at ($(constzero)!0.5!(agree)$)
    {one-sided\\limits agree};
  \node [barlabelB] at ($(constzero)!0.5!(differ)$)
    {one-sided\\limits disagree};
  \node [barlabelB] at ($(conttest)!0.5!(yes)$) {yes};
  \node [barlabelB] at ($(conttest)!0.5!(no)$) {no};
  \node [barlabelB] at ($(root)!0.5!(dirsub)$) {calculations uses\\limit laws};
  % Rides the bar into Squeeze the way the other decisions ride theirs.
  \node [barlabelB] at ($(dirsub)!0.5!(squeeze)$) {recognize\\inequalities};
  % Graphs lifted from plot-limit-intro.pdf and parked beside the node each
  % one illustrates.  \includegraphics reads the built PDF, so the path is
  % relative to this directory -- and plot-limit-intro.pdf has to exist before
  % this file is compiled.
  % Anchored off the node each one illustrates, so they follow if the layout
  % moves.  `annotation' forces a 3.5cm ragged text box and these images are
  % 3.81cm wide, so reset the text width to let each node size to its image --
  % otherwise every one of them overfulls and hangs past its own node box.
  \node [concept color=white, annotation, text width=, above]
    at (twosided.north) {\includegraphics[width=0.75in, page=3]{plot-limit-intro}};
  \node [concept color=white, annotation, text width=, right]
    at (leftlim.east) {\includegraphics[width=0.75in, page=1]{plot-limit-intro}};
  \node [concept color=white, annotation, text width=, above]
    at (rightlim.north) {\includegraphics[width=0.75in, page=2]{plot-limit-intro}};

  % The one-sided continuity pictures, smaller because there are four of them
  % in a crowded corner.  The last is stacked above its sibling rather than
  % anchored straight off the node: everything below rightcont is taken.
  \node [concept color=white, annotation, text width=, left]
    at (leftcont.west) {\includegraphics[width=0.75in, page=13]{plot-limit-intro}};
  \node [concept color=white, annotation, text width=, below left]
    at (leftcont.south) {\includegraphics[width=0.75in, page=14]{plot-limit-intro}};
  \node [concept color=white, annotation, text width=, below]
    at (rightcont.south) {\includegraphics[width=0.75in, page=15]{plot-limit-intro}};
  \node [concept color=white, annotation, text width=, below]
    at ([shift={(-2.1cm,0)}]rightcont.south) {\includegraphics[width=0.75in, page=16]{plot-limit-intro}};



  % Nudged clear of the connection it labels -- the label sits on the main
  % layer and the bar below it is saturated, so text straight on top of it
  % would not read.
  % The arrow is vertical on this page, so the label runs up it.  rotate=90
  % turns the node's own south edge to face east, so anchor=south still
  % sets the text to the left of the bar.
  \node [barlabelB, rotate=90, anchor=south]
    at ($(conttest)!0.5!(dirsub)$) {testing for continuity\\uses these techniques};

  % The two decision trees reach the same place by different routes.
  \begin{pgfonlayer}{background}
    \draw [concept connection, substline, ->, >=stealth, shorten >=4pt]
      (algno) edge (yes);
    % Directed: continuity is what borrows the techniques, not the other way
    % round.  Shortened so the head clears the concept's blob outline.  Drawn
    % in the root's own grey, so it reads as part of the map's spine rather
    % than as one of the coloured cross-links.
    \draw [concept connection, gray!20, ->, >=stealth, shorten >=4pt]
      (conttest) edge (dirsub);
    \draw [concept connection, supp!60, ->, >=stealth, shorten >=4pt]
      (cov) edge (known);
    % A circular arc, bowed away from the root so it threads the gap between
    % the root and its branches.  No heads, to match the other plain link.
    \draw [concept connection, pink]
      (constzero.east) to [bend left=20] (infinite);
  \end{pgfonlayer}
\end{tikzpicture} 
}
\end{mmpage}

\end{document}
